% ---------------------------------------------------------------------------
% Study report · The unified dataset, and what imputation can do for it
%
% Build:  pdflatex unified_imputation_report.tex  (twice, for the ToC)
% Figures are read from ../outputs/ and are produced by the paired notebooks.
% ---------------------------------------------------------------------------
\documentclass[11pt,a4paper]{article}

\newcommand{\runninghead}{Unified dataset and imputation · 2018--2026}

% ---------------------------------------------------------------------------
% reportstyle.tex — shared preamble for the study reports under studies/.
%
% Kept deliberately inside the package set shipped with TeX Live "basic":
% no tcolorbox, no siunitx, no enumitem, no titlesec. Everything below
% compiles with a stock pdflatex installation.
%
% Usage, from a report directory one level below a study folder:
%     % ---------------------------------------------------------------------------
% reportstyle.tex — shared preamble for the study reports under studies/.
%
% Kept deliberately inside the package set shipped with TeX Live "basic":
% no tcolorbox, no siunitx, no enumitem, no titlesec. Everything below
% compiles with a stock pdflatex installation.
%
% Usage, from a report directory one level below a study folder:
%     % ---------------------------------------------------------------------------
% reportstyle.tex — shared preamble for the study reports under studies/.
%
% Kept deliberately inside the package set shipped with TeX Live "basic":
% no tcolorbox, no siunitx, no enumitem, no titlesec. Everything below
% compiles with a stock pdflatex installation.
%
% Usage, from a report directory one level below a study folder:
%     \input{../../_shared/reportstyle.tex}
%     \graphicspath{{../outputs/}}
% ---------------------------------------------------------------------------

\usepackage[T1]{fontenc}
\usepackage[utf8]{inputenc}
\usepackage{textcomp}
\usepackage{lmodern}
\usepackage{microtype}
\usepackage[margin=2.4cm,headheight=15pt]{geometry}
\usepackage{graphicx}
\usepackage{booktabs}
\usepackage{tabularx}
\usepackage{array}
\usepackage{amsmath}
\usepackage{xcolor}
\usepackage[font=small,labelfont=bf,justification=justified]{caption}
\usepackage{float}
\usepackage{fancyhdr}
\usepackage[hidelinks]{hyperref}

% --- Palette ---------------------------------------------------------------
\definecolor{rulegrey}{HTML}{B9C0C7}
\definecolor{fillgrey}{HTML}{F2F4F6}
\definecolor{failred}{HTML}{A5202B}
\definecolor{passgreen}{HTML}{1B6B3A}
\definecolor{accent}{HTML}{1F4E79}

% --- Semantic shorthands ---------------------------------------------------
\newcommand{\degC}{\textdegree C}
\newcommand{\mdegC}{mdeg/\textdegree C}
\newcommand{\fail}{\textcolor{failred}{\textbf{fail}}}
\newcommand{\pass}{\textcolor{passgreen}{\textbf{pass}}}
\newcommand{\worse}{\textcolor{failred}{worse}}
\newcommand{\better}{\textcolor{passgreen}{better}}

% --- Highlighted panel -----------------------------------------------------
% #1 = heading, #2 = body
\newcommand{\panel}[2]{%
  \begin{center}%
  \setlength{\fboxsep}{9pt}%
  \fcolorbox{rulegrey}{fillgrey}{%
    \begin{minipage}{0.93\textwidth}%
      {\bfseries\color{accent}#1}\par\medskip
      #2
    \end{minipage}}%
  \end{center}}

% --- Numeric column: right aligned, fixed width ----------------------------
\newcolumntype{R}[1]{>{\raggedleft\arraybackslash}p{#1}}
\newcolumntype{L}[1]{>{\raggedright\arraybackslash}p{#1}}
\newcolumntype{Y}{>{\raggedright\arraybackslash}X}

% --- Table look ------------------------------------------------------------
\renewcommand{\arraystretch}{1.15}
\setlength{\tabcolsep}{6pt}

% --- Headers ---------------------------------------------------------------
\pagestyle{fancy}
\fancyhf{}
\fancyhead[L]{\small\itshape\runninghead}
\fancyhead[R]{\small\thepage}
\renewcommand{\headrulewidth}{0.4pt}

% --- Section spacing without titlesec --------------------------------------
\makeatletter
\renewcommand\section{\@startsection{section}{1}{\z@}%
  {-3.2ex \@plus -1ex \@minus -.2ex}{2.0ex \@plus.2ex}%
  {\normalfont\Large\bfseries\color{accent}}}
\renewcommand\subsection{\@startsection{subsection}{2}{\z@}%
  {-2.6ex \@plus -1ex \@minus -.2ex}{1.4ex \@plus.2ex}%
  {\normalfont\large\bfseries}}
\makeatother

% --- Figure defaults -------------------------------------------------------
\setkeys{Gin}{width=\linewidth}
\captionsetup[figure]{skip=6pt}

%     \graphicspath{{../outputs/}}
% ---------------------------------------------------------------------------

\usepackage[T1]{fontenc}
\usepackage[utf8]{inputenc}
\usepackage{textcomp}
\usepackage{lmodern}
\usepackage{microtype}
\usepackage[margin=2.4cm,headheight=15pt]{geometry}
\usepackage{graphicx}
\usepackage{booktabs}
\usepackage{tabularx}
\usepackage{array}
\usepackage{amsmath}
\usepackage{xcolor}
\usepackage[font=small,labelfont=bf,justification=justified]{caption}
\usepackage{float}
\usepackage{fancyhdr}
\usepackage[hidelinks]{hyperref}

% --- Palette ---------------------------------------------------------------
\definecolor{rulegrey}{HTML}{B9C0C7}
\definecolor{fillgrey}{HTML}{F2F4F6}
\definecolor{failred}{HTML}{A5202B}
\definecolor{passgreen}{HTML}{1B6B3A}
\definecolor{accent}{HTML}{1F4E79}

% --- Semantic shorthands ---------------------------------------------------
\newcommand{\degC}{\textdegree C}
\newcommand{\mdegC}{mdeg/\textdegree C}
\newcommand{\fail}{\textcolor{failred}{\textbf{fail}}}
\newcommand{\pass}{\textcolor{passgreen}{\textbf{pass}}}
\newcommand{\worse}{\textcolor{failred}{worse}}
\newcommand{\better}{\textcolor{passgreen}{better}}

% --- Highlighted panel -----------------------------------------------------
% #1 = heading, #2 = body
\newcommand{\panel}[2]{%
  \begin{center}%
  \setlength{\fboxsep}{9pt}%
  \fcolorbox{rulegrey}{fillgrey}{%
    \begin{minipage}{0.93\textwidth}%
      {\bfseries\color{accent}#1}\par\medskip
      #2
    \end{minipage}}%
  \end{center}}

% --- Numeric column: right aligned, fixed width ----------------------------
\newcolumntype{R}[1]{>{\raggedleft\arraybackslash}p{#1}}
\newcolumntype{L}[1]{>{\raggedright\arraybackslash}p{#1}}
\newcolumntype{Y}{>{\raggedright\arraybackslash}X}

% --- Table look ------------------------------------------------------------
\renewcommand{\arraystretch}{1.15}
\setlength{\tabcolsep}{6pt}

% --- Headers ---------------------------------------------------------------
\pagestyle{fancy}
\fancyhf{}
\fancyhead[L]{\small\itshape\runninghead}
\fancyhead[R]{\small\thepage}
\renewcommand{\headrulewidth}{0.4pt}

% --- Section spacing without titlesec --------------------------------------
\makeatletter
\renewcommand\section{\@startsection{section}{1}{\z@}%
  {-3.2ex \@plus -1ex \@minus -.2ex}{2.0ex \@plus.2ex}%
  {\normalfont\Large\bfseries\color{accent}}}
\renewcommand\subsection{\@startsection{subsection}{2}{\z@}%
  {-2.6ex \@plus -1ex \@minus -.2ex}{1.4ex \@plus.2ex}%
  {\normalfont\large\bfseries}}
\makeatother

% --- Figure defaults -------------------------------------------------------
\setkeys{Gin}{width=\linewidth}
\captionsetup[figure]{skip=6pt}

%     \graphicspath{{../outputs/}}
% ---------------------------------------------------------------------------

\usepackage[T1]{fontenc}
\usepackage[utf8]{inputenc}
\usepackage{textcomp}
\usepackage{lmodern}
\usepackage{microtype}
\usepackage[margin=2.4cm,headheight=15pt]{geometry}
\usepackage{graphicx}
\usepackage{booktabs}
\usepackage{tabularx}
\usepackage{array}
\usepackage{amsmath}
\usepackage{xcolor}
\usepackage[font=small,labelfont=bf,justification=justified]{caption}
\usepackage{float}
\usepackage{fancyhdr}
\usepackage[hidelinks]{hyperref}

% --- Palette ---------------------------------------------------------------
\definecolor{rulegrey}{HTML}{B9C0C7}
\definecolor{fillgrey}{HTML}{F2F4F6}
\definecolor{failred}{HTML}{A5202B}
\definecolor{passgreen}{HTML}{1B6B3A}
\definecolor{accent}{HTML}{1F4E79}

% --- Semantic shorthands ---------------------------------------------------
\newcommand{\degC}{\textdegree C}
\newcommand{\mdegC}{mdeg/\textdegree C}
\newcommand{\fail}{\textcolor{failred}{\textbf{fail}}}
\newcommand{\pass}{\textcolor{passgreen}{\textbf{pass}}}
\newcommand{\worse}{\textcolor{failred}{worse}}
\newcommand{\better}{\textcolor{passgreen}{better}}

% --- Highlighted panel -----------------------------------------------------
% #1 = heading, #2 = body
\newcommand{\panel}[2]{%
  \begin{center}%
  \setlength{\fboxsep}{9pt}%
  \fcolorbox{rulegrey}{fillgrey}{%
    \begin{minipage}{0.93\textwidth}%
      {\bfseries\color{accent}#1}\par\medskip
      #2
    \end{minipage}}%
  \end{center}}

% --- Numeric column: right aligned, fixed width ----------------------------
\newcolumntype{R}[1]{>{\raggedleft\arraybackslash}p{#1}}
\newcolumntype{L}[1]{>{\raggedright\arraybackslash}p{#1}}
\newcolumntype{Y}{>{\raggedright\arraybackslash}X}

% --- Table look ------------------------------------------------------------
\renewcommand{\arraystretch}{1.15}
\setlength{\tabcolsep}{6pt}

% --- Headers ---------------------------------------------------------------
\pagestyle{fancy}
\fancyhf{}
\fancyhead[L]{\small\itshape\runninghead}
\fancyhead[R]{\small\thepage}
\renewcommand{\headrulewidth}{0.4pt}

% --- Section spacing without titlesec --------------------------------------
\makeatletter
\renewcommand\section{\@startsection{section}{1}{\z@}%
  {-3.2ex \@plus -1ex \@minus -.2ex}{2.0ex \@plus.2ex}%
  {\normalfont\Large\bfseries\color{accent}}}
\renewcommand\subsection{\@startsection{subsection}{2}{\z@}%
  {-2.6ex \@plus -1ex \@minus -.2ex}{1.4ex \@plus.2ex}%
  {\normalfont\large\bfseries}}
\makeatother

% --- Figure defaults -------------------------------------------------------
\setkeys{Gin}{width=\linewidth}
\captionsetup[figure]{skip=6pt}

\graphicspath{{../outputs/}}

\title{\vspace{-1.2cm}\textbf{One dataset for the whole archive, and what
imputation can do for it}\\[2mm]
\large Every station, both instrument eras, 26 July 2018 to 13 August 2026}
\author{Study report · \texttt{studies/unified\_dataset/}}
\date{13 August 2026}

\begin{document}
\maketitle
\thispagestyle{fancy}

\begin{abstract}
\noindent
Every earlier study in this project began by reading the raw archive and ended
by discarding most of it. This one builds the dataset those studies should have
started from --- all three legacy stations, both column eras, cleaned,
compensated and carrying provenance on every row --- and then asks what can be
done about the holes.

\textbf{The archive is emptier than its span suggests.} Eight years of calendar
time contain \textbf{2061 days} of inclinometer readings at the target station,
70\,\% of the span; the other two stations hold 1119 and 984 days and both
stopped recording years before the February 2025 instrument replacement. The
network narrowed from three locations to one in stages, not at a single moment.

\textbf{The missingness is the worst kind available.} On every gap, every
channel is missing at once: air temperature, humidity, battery and inclinometer
fail together on \textbf{100\,\%} of the missing hours. This is outage-driven
missingness, and it has a consequence that governs the whole study --- no method
conditioned on a covariate can ever fill a real gap here, because the covariates
are gone with the target.

\textbf{The benchmark cannot see that, and the study reports the discrepancy.}
Injected gaps have their drivers intact, so a driver-based filler scores best at
almost every length and would be useless in production. This is the
exchangeability assumption behind any conformal guarantee, failing concretely
and measurably.

\textbf{The two eras are two instruments, and a constant does not reconcile
them.} The step at the changeover is $-132.5$\,mdeg once the season is removed,
but the diurnal amplitude rises from 13.4 to 17.9\,mdeg across it --- a
\textbf{ratio of 1.34}. The instruments differ in gain as well as in zero, so
the joined series is offered as a clearly named derived column and never as the
default.

\textbf{The record does not have to stay in pieces.} Without filling anything,
the longest strictly continuous stretch in eight years is \textbf{108 days}.
Imputation raises that to \textbf{988.9 days --- 2.71 annual cycles} --- at
14.5\,\% invented, or to \textbf{617.2 days} at 3.2\,\% if the stretch is
required to stay within one instrument era. Both carry a calibrated interval on
every invented value, and past a 15\,mdeg tolerance no further record can be
bought at any price.
\end{abstract}

\tableofcontents

% ===========================================================================
\section{What this study builds, and why}
\label{sec:intro}

\subsection{Two deliverables}

The work is split into two notebooks, deliberately.

\texttt{build\_unified\_dataset.py} is the \textbf{ingest}, and it is the only place
in the project that reads the raw archive. It is meant to be re-run whenever the
archive grows: it discovers the extent itself, so no date has to be edited. It
writes one table per station to \texttt{data/interim/unified/}, plus a manifest
recording how the build was made.

\texttt{imputation\_study.py} is this report. It consumes that output and never
touches an \texttt{.adc} file.

The separation is not tidiness. The ingest writes only measurements, and every
row it writes carries \texttt{inc\_source = observed}. Imputation happens on the
far side of that boundary, so a filled value can never be mistaken for a reading
by a notebook three steps downstream.

\subsection{Three rules the dataset obeys}

\panel{The rules, and where they come from}{%
\textbf{The two eras are two instruments.} The package installed on 2025-02-21
replaced the legacy network rather than augmenting it, and its inclinometer
shares no baseline with legacy block b2 (\texttt{docs/raw-data-format.md}
Section 3.3). The primary \texttt{inc} and \texttt{inc\_comp} columns keep the
baselines apart; \texttt{era} and \texttt{instrument} say which is which.

\medskip
\textbf{The absolute level means nothing.} Only changes carry structural
information (Section 7.1 of the same document). Every level here is referred to
an arbitrary anchor.

\medskip
\textbf{A filled value is not a measurement.} Provenance is carried on every
row, through every product, to the end.}

\subsection{What the dataset contains}

One CSV per station on the hourly grid, spanning the whole archive.

\begin{center}
\small
\begin{tabularx}{\textwidth}{L{3.4cm}Y}
\toprule
\textbf{Column} & \textbf{Meaning} \\
\midrule
\texttt{batt}, \texttt{tair}, \texttt{rh} & Supply voltage, air temperature,
relative humidity \\
\texttt{inc} & Inclinometer, millidegrees about the instrument's calibrated
zero. Its absolute level carries no structural information \\
\texttt{sr}, \texttt{twall} & Solar radiation and wall temperature. Current era,
target station only \\
\texttt{inc\_comp} & Compensated at the documented coefficient, anchored
\emph{within each era separately} \\
\texttt{inc\_comp\_joined} & The same with the estimated instrument offset
removed, so the eras form one series. Target station only. Carries an estimated
constant \\
\texttt{era}, \texttt{instrument} & Which era, and which physical instrument \\
\texttt{inc\_source} & \texttt{observed} or \texttt{missing}. The ingest fills
nothing \\
\bottomrule
\end{tabularx}
\end{center}

The compensation uses the documented coefficient of 5\,mdeg per \degC. That
value has never been calibrated for these instruments, and the legacy study
measured its influence on a headline error at $-58$\,\% to $+183$\,\%. It is
used here for continuity with every earlier study, and the reader should treat
absolute levels accordingly.

% ===========================================================================
\section{What is actually there}
\label{sec:present}

\subsection{Three stations, three different records}

\begin{center}
\small
\begin{tabular}{lllR{1.6cm}R{1.7cm}R{1.6cm}}
\toprule
\textbf{Station} & \textbf{Eras} & \textbf{Last reading} &
\textbf{Span [d]} & \textbf{Observed [d]} & \textbf{of archive} \\
\midrule
\texttt{st01} & legacy         & 2022-06-01 & 1406.3 & 1118.7 & 38.0\,\% \\
\texttt{st02} & legacy+current & 2026-08-13 & 2940.0 & 2061.0 & 70.1\,\% \\
\texttt{st03} & legacy         & 2023-06-24 & 1794.6 &  984.1 & 33.5\,\% \\
\bottomrule
\end{tabular}
\end{center}

\panel{The network failed in stages, not at the changeover}{%
The data-quality report records that the current-era package ``replaced'' the
legacy network, which is true of the file format and misleading about the
instrumentation. \textbf{Station 01 stopped recording in June 2022 and station
03 in June 2023} --- two years and eight months before the replacement
respectively.

\medskip
For most of the archive there was never a three-station network to compare
across. Any cross-station test --- including the one most likely to settle the
sign contradiction documented in the two prediction studies --- is confined to
the period before June 2022.}

\begin{figure}[H]
\includegraphics{U_F03_station_completeness}
\caption{Daily availability of the inclinometer at each station. The staggered
right-hand edges are the point: the three records end three years apart.}
\end{figure}

\subsection{The target station, era by era}

\begin{center}
\small
\begin{tabular}{lllR{1.5cm}R{1.5cm}R{1.3cm}R{1.3cm}}
\toprule
\textbf{Era} & \textbf{From} & \textbf{To} & \textbf{Span [d]} &
\textbf{Cycles} & \textbf{\texttt{inc}} & \textbf{\texttt{twall}} \\
\midrule
legacy  & 2018-07-26 & 2025-02-20 & 2401.0 & 6.57 & 68.5\,\% & --- \\
current & 2025-02-21 & 2026-08-13 &  538.0 & 1.47 & 77.2\,\% & 36.8\,\% \\
\bottomrule
\end{tabular}
\end{center}

The current era is the better-instrumented and the better-covered of the two,
and it is short. Wall temperature is present on barely a third of it, because
the probe-failure sentinel and the acquisition outages interleave.

\begin{figure}[H]
\includegraphics{U_F04_channel_availability}
\caption{Daily availability of every channel at the target station. Solar
radiation and wall temperature begin only at the changeover; wall temperature is
then interrupted repeatedly by its own failure sentinel.}
\end{figure}

\subsection{The calendar}

\begin{figure}[H]
\includegraphics{U_F05_calendar}
\caption{Fraction of each day observed, by year and day of year. Outages appear
as horizontal bands and are not confined to any season.}
\end{figure}

% ===========================================================================
\section{The instrument changeover}
\label{sec:changeover}

The two eras are adjacent --- the legacy record ends on 2025-02-19 and the
current one begins on 2025-02-21, an interruption of about a day --- so a level
comparison across the boundary is possible. It is also misleading if taken
naively, because the two comparison windows sit at different points of the
annual cycle.

\begin{center}
\small
\begin{tabular}{lR{2.6cm}}
\toprule
\textbf{Quantity} & \textbf{Value} \\
\midrule
Interruption                       & 1.04 d \\
Raw step                           & $-127.67$ mdeg \\
Attributable to the season         & $+4.81$ mdeg \\
\textbf{Attributed to the instrument} & $\mathbf{-132.48}$ \textbf{mdeg} \\
\midrule
Diurnal amplitude, legacy side     & 13.38 mdeg \\
Diurnal amplitude, current side    & 17.91 mdeg \\
\textbf{Gain ratio}                & $\mathbf{1.34}$ \\
\bottomrule
\end{tabular}
\end{center}

\panel{A constant offset does not reconcile the two instruments}{%
The offset is estimable and the estimate is reported. The diagnostic beside it
is the one that matters: the diurnal amplitude --- the signal's response to the
daily thermal cycle, which is a property of the wall and not of the logger ---
rises by \textbf{34\,\%} across the changeover.

\medskip
The wall did not become a third more responsive to sunlight overnight. Either
the new instrument has a different gain, or it is mounted where the thermal
response is larger, or both. A single additive constant cannot absorb a
multiplicative difference, so \texttt{inc\_comp\_joined} reconciles the
\emph{levels} of the two eras and does not reconcile their \emph{scales}.

\medskip
This is why the joined series is a derived column and not the default, and why
any analysis spanning the changeover should either work in the difference domain
or report results per era.}

\begin{figure}[H]
\includegraphics{U_F02_era_join}
\caption{The compensated series before and after the join. The offset removes
the level discontinuity; it cannot remove the amplitude difference.}
\end{figure}

% ===========================================================================
\section{What is missing, and why}
\label{sec:missing}

\subsection{The anatomy of the gaps}

\begin{center}
\small
\begin{tabular}{lR{1.4cm}R{1.7cm}R{1.8cm}R{1.5cm}}
\toprule
\textbf{Band} & \textbf{Gaps} & \textbf{Hours lost} &
\textbf{Share} & \textbf{Median} \\
\midrule
$\leq$ 3 h & 426 &   619 &  2.9\,\% &  1 h \\
3--6 h     &  60 &   284 &  1.3\,\% &  4 h \\
6--24 h    &  72 &   906 &  4.3\,\% & 11 h \\
1--3 d     &   2 &    65 &  0.3\,\% & 33 h \\
3--7 d     &   0 &     0 &  0.0\,\% & --- \\
7--30 d    &   3 & 1\,415 &  6.7\,\% & 491 h \\
\textbf{$>$ 30 d} & \textbf{6} & \textbf{17\,809} &
\textbf{84.4\,\%} & 1\,396 h \\
\bottomrule
\end{tabular}
\end{center}

569 gaps hold 21\,098 missing hours, and the distribution is extremely uneven.
\textbf{The 560 gaps of three days or less account for 8.9\,\% of the loss; six
gaps longer than thirty days account for 84.4\,\%.} Almost everything that is
missing is missing in a handful of long outages.

\begin{figure}[H]
\includegraphics{U_F06_gap_anatomy}
\caption{Left: share of missing hours by length band, annotated with the number
of gaps. Centre: hours lost by season of onset. Right: every gap, by date and
length, on a logarithmic scale.}
\end{figure}

\subsection{The nine that matter}

\begin{center}
\small
\begin{tabular}{llR{1.5cm}l}
\toprule
\textbf{From} & \textbf{To} & \textbf{Days} & \textbf{Season} \\
\midrule
\textbf{2022-04-09} & \textbf{2023-06-21} & \textbf{437.6} & MAM \\
2026-03-06 & 2026-06-18 & 103.9 & MAM \\
2018-09-23 & 2018-11-21 &  59.0 & SON \\
2020-06-04 & 2020-07-31 &  57.3 & JJA \\
2024-08-27 & 2024-10-09 &  42.7 & JJA \\
2024-05-04 & 2024-06-14 &  41.5 & MAM \\
2019-01-14 & 2019-02-04 &  20.8 & DJF \\
2019-04-20 & 2019-05-11 &  20.5 & MAM \\
2025-09-25 & 2025-10-13 &  17.8 & SON \\
\bottomrule
\end{tabular}
\end{center}

These nine cut the eight-year archive into ten pieces. Their onsets are spread
across all four seasons, so there is no seasonal maintenance pattern to exploit.

\subsection{Why the record is missing}

\panel{Every channel fails at once}{%
Of the hours in which the inclinometer is missing, the fraction in which each
other channel is \emph{also} missing:

\medskip
\begin{tabular}{lR{2.0cm}}
\texttt{batt}  & 100.0\,\% \\
\texttt{tair}  & 100.0\,\% \\
\texttt{rh}    & 100.0\,\% \\
\texttt{sr}    & 100.0\,\% \\
\texttt{twall} & 100.0\,\% \\
\end{tabular}

\medskip
There is no covariate that is ever observed while the target is absent. The
missingness is not random, and it is not conditional on something the record
still contains --- it is an \textbf{outage}, and an outage takes the whole
station with it.

\medskip
\textbf{The consequence is hard and it governs Section~\ref{sec:fill}: no method
conditioned on a covariate can fill a real gap in this record.} Not because such
methods are poor, but because their inputs are missing exactly when they would
be needed.}

This also settles a question the framework in the companion manuscript raises
but cannot answer from the extended-sensor era alone. Its regressor-diagnosis
stage applies an availability filter --- does a candidate exist during the
intervals that must be reconstructed? --- and warns that on-structure channels
sharing the station's power and communications will fail it by construction.
Here that filter does not merely reduce the candidate pool. \textbf{It empties
it.}

% ===========================================================================
\section{What can be filled, and how well}
\label{sec:fill}

Six fillers of increasing ambition, measured against gaps of known length
injected where the record is complete. No filler sees the values it is asked to
reproduce. The range matters because no single method is right at every length:
interpolation is unbeatable across an hour and meaningless across a month, while
a seasonal model knows nothing useful about an hour and is most of what remains
after a week.

\begin{center}
\small
\begin{tabularx}{\textwidth}{L{3.0cm}Y}
\toprule
\textbf{Filler} & \textbf{What it knows} \\
\midrule
\texttt{interpolate}    & The two flanking observations, and nothing else \\
\texttt{seasonal\_naive} & That there is a daily cycle \\
\texttt{kalman}          & A local linear trend plus a daily seasonal, in state-space form. The only filler with a native notion of its own uncertainty \\
\texttt{harmonic}        & The annual and diurnal cycles, fitted by least squares. Its error does not grow with gap length \\
\texttt{drivers}         & A ridge regression on air temperature and humidity \\
\texttt{neuralprophet}   & Trend and seasonality from the project's own decomposition \\
\bottomrule
\end{tabularx}
\end{center}

\subsection{Error against gap length}

\begin{center}
\small
\begin{tabular}{rR{1.6cm}R{1.8cm}R{1.4cm}R{1.8cm}R{1.4cm}}
\toprule
\textbf{Gap} & \textbf{interp.} & \textbf{seas.\ naive} &
\textbf{kalman} & \textbf{harmonic} & \textbf{drivers} \\
\midrule
1 h     & \textbf{1.22} &  4.58 &  \textbf{1.04} &  3.84 &  1.83 \\
3 h     &  1.61 &  5.67 &  1.75 &  3.73 & \textit{1.14} \\
6 h     &  \textbf{2.79} &  4.52 &  2.99 &  4.00 & \textit{2.48} \\
12 h    &  8.75 &  9.04 &  \textbf{8.55} &  9.98 & \textit{7.25} \\
24 h    &  9.83 &  \textbf{5.12} &  5.71 &  6.13 & \textit{4.00} \\
48 h    &  9.54 &  4.94 &  \textbf{4.18} &  4.80 & \textit{2.07} \\
72 h    & 14.70 &  \textbf{7.31} &  \textbf{7.31} &  8.24 & \textit{5.50} \\
168 h   & 10.68 &  \textbf{8.05} &  9.86 & 10.02 & \textit{6.98} \\
336 h   & 21.79 & \textbf{11.50} & 13.07 & 12.84 & \textit{8.78} \\
720 h   & 15.79 & \textbf{13.62} & 16.69 & 15.56 & 14.49 \\
1080 h  & 16.13 & 13.68 & \textbf{11.29} & 14.63 & \textit{10.65} \\
\bottomrule
\end{tabular}
\end{center}

\noindent{\footnotesize Mean absolute error in mdeg, against a signal standard
deviation of about 52\,mdeg. \textbf{Bold}: best method that can actually run on
a real gap. \textit{Italic}: the driver filler, which cannot --- see
Section~\ref{sec:illusion}.}

\begin{figure}[H]
\includegraphics{U_F07_filler_benchmark}
\caption{Left: mean absolute error against gap length. Right: the width of the
calibrated 90\,\% interval for the same method and length.}
\end{figure}

\panel{No method wins everywhere, and the crossover is sharp}{%
\textbf{Below six hours, interpolation and the Kalman smoother are
interchangeable and both are excellent} --- around 1 to 3\,mdeg, which is
2 to 6\,\% of the signal's standard deviation. Nothing more elaborate helps,
because over a few hours the series really is close to a straight line between
its endpoints.

\medskip
\textbf{Past twelve hours interpolation collapses and never recovers.} At 24
hours it is at 9.8\,mdeg while the seasonal-naive filler is at 5.1, because a
straight line drawn across a day removes the diurnal cycle the data actually
contains. This is the single clearest result in the benchmark and it sets the
six-hour boundary used in every policy below.

\medskip
\textbf{Beyond a day, the choice is between knowing the daily cycle and knowing
the local state}, and they trade places repeatedly: seasonal-naive wins at 24,
72, 168, 336 and 720 hours, the Kalman smoother at 48 and 1080. The differences
are small compared with the error itself, so the policy's choice among them
matters far less than the decision to stop interpolating.}

\subsection{The uncertainty attached to a filled value}

A mean absolute error says how wrong a filler is on average. It does not say how
wrong a particular value might be, which is what a consumer needs in order to
weight or exclude it. The per-trial errors are therefore used directly as a
calibration sample: the interval covering 90\,\% of them is, by construction,
the interval that would have covered 90\,\% of the errors observed. No
distributional assumption is made --- these errors are neither Gaussian nor
independent, and assuming either would understate them.

\begin{center}
\small
\begin{tabular}{rR{1.7cm}R{1.9cm}R{1.4cm}R{1.8cm}}
\toprule
\textbf{Gap} & \textbf{interp.} & \textbf{seas.\ naive} &
\textbf{kalman} & \textbf{harmonic} \\
\midrule
1 h    &  3.16 &  8.39 &  2.06 &  7.63 \\
6 h    &  6.20 & 10.17 &  5.81 &  8.95 \\
24 h   & 22.71 &  9.56 &  8.51 &  9.00 \\
168 h  & 21.19 & 13.92 & 15.75 & 14.95 \\
720 h  & 29.20 & 21.44 & 29.31 & 25.18 \\
1080 h & 31.62 & 26.33 & 19.94 & 28.54 \\
\bottomrule
\end{tabular}
\end{center}

\noindent{\footnotesize Half-width of the 90\,\% interval, in mdeg.}

The interval widens with gap length for every method, as it must. It is also
consistently two to three times the mean absolute error, which is the honest
consequence of a long-tailed error distribution: the typical filled value is
good, and the occasional one is not.

\subsection{The driver filler, and why its score is not to be believed}
\label{sec:illusion}

The driver filler has the lowest error at eight of the eleven lengths tested. It
is also the one method in the set that \textbf{cannot be used at all}.

Section~\ref{sec:missing} established that every channel of this station fails
together: on 100\,\% of the missing hours, air temperature and humidity are
missing too. The benchmark cannot see this, because it injects its gaps where
the record is complete --- so the drivers are present throughout every injected
gap, and absent throughout every real one.

\panel{This is the exchangeability assumption, failing measurably}{%
Any conformal guarantee holds under exchangeability between the calibration
sample and the deployment conditions. Here the two differ in a way that is not
subtle and not statistical: the calibration gaps have their predictors, the
deployment gaps do not.

\medskip
Had the policy been allowed to select on benchmark score alone, it would have
chosen the driver filler at eight lengths, and the resulting fill would have
degraded silently to interpolation while carrying the driver filler's error bars
--- claiming 2.07\,mdeg at 48 hours where interpolation delivers 9.54.

\medskip
The study therefore measures the method and then refuses it. What that refusal
costs is small and is stated: between 0.0 and 2.7\,mdeg of accuracy per band,
1.05\,mdeg on average.}

This is the concrete case for the availability criterion that the companion
manuscript's regressor-diagnosis stage applies. At this site the criterion does
not merely reorder the candidates --- it removes every one of them.

% ===========================================================================
\section{The longest stretch}
\label{sec:stretch}

Every extra hour of continuous record costs invented data. The frontier is the
exchange rate, and it is the answer this study exists to give.

\subsection{Without filling anything}

\panel{Eight years of archive contain 108 consecutive days}{%
The longest run in which no hour is missing, over the whole record, runs from
\textbf{2025-05-05 to 2025-08-21} --- \textbf{108.2 days}, or 0.30 of an annual
cycle. It sits in the current era, which is the shortest and best-covered part
of the archive.

\medskip
That is what the record offers to any method requiring strict continuity, and it
is why every earlier study in this project tolerated gaps rather than demanding
their absence.}

\subsection{The frontier}

\begin{center}
\small
\begin{tabular}{lR{1.6cm}R{1.4cm}llR{2.0cm}}
\toprule
\textbf{Tolerance} & \textbf{Stretch} & \textbf{Cycles} & \textbf{From} &
\textbf{To} & \textbf{Fabricated} \\
\midrule
--- (observed) & 108.2 d & 0.30 & 2025-05-05 & 2025-08-21 & 0.00\,\% \\
2 mdeg   & 169.5 d & 0.46 & 2023-10-30 & 2024-04-16 &  0.61\,\% \\
5 mdeg   & 211.5 d & 0.58 & 2021-08-20 & 2022-03-19 &  1.60\,\% \\
10 mdeg  & 617.2 d & 1.69 & 2020-07-31 & 2022-04-09 &  3.20\,\% \\
\textbf{15 mdeg} & \textbf{988.9 d} & \textbf{2.71} &
\textbf{2023-06-21} & \textbf{2026-03-06} & \textbf{14.52\,\%} \\
20 mdeg  & 988.9 d & 2.71 & 2023-06-21 & 2026-03-06 & 14.52\,\% \\
40 mdeg  & 988.9 d & 2.71 & 2023-06-21 & 2026-03-06 & 14.52\,\% \\
\bottomrule
\end{tabular}
\end{center}

\noindent{\footnotesize Fabricated share is measured \emph{inside} the winning
stretch, not across the record.}

\begin{figure}[H]
\includegraphics{U_F08_stretch_frontier}
\caption{Left: longest continuous stretch per policy, annotated with the annual
cycles it contains; the dashed line is what the record offers unfilled. Right:
the same stretches against the fabricated share paid for them.}
\end{figure}

\panel{The frontier has one sharp step and then stops}{%
Between 2 and 10\,mdeg of tolerance the stretch grows steadily, and everything
gained is cheap: at 10\,mdeg the record yields \textbf{617 days for 3.2\,\%
fabrication}, which is 5.7 times what strict continuity offers.

\medskip
At 15\,mdeg the stretch jumps to \textbf{988.9 days}, because that is the point
at which the two 41-- and 43-day outages of 2024 become fillable and three
segments merge into one. The jump costs a step change in fabrication, from
3.2\,\% to \textbf{14.5\,\%}.

\medskip
Past 15\,mdeg nothing further is bought at any price. The remaining boundaries
are the 59-day, 104-day and 438-day outages, all longer than the longest gap the
benchmark can calibrate, and the 438-day one is longer than most of the segments
it separates. \textbf{Raising the tolerance to 40\,mdeg --- most of the signal's
standard deviation --- buys not one additional day.}}

\subsection{What the era join is worth}

The 988.9-day stretch runs from June 2023 to March 2026 and therefore
\textbf{crosses the February 2025 instrument changeover}. It exists only in the
joined column, and so rests on the estimated offset of
Section~\ref{sec:changeover}.

Confining every stretch to a single era answers what is available without that
estimate.

\begin{center}
\small
\begin{tabular}{lR{2.0cm}R{1.6cm}R{2.4cm}}
\toprule
& \textbf{Stretch} & \textbf{Cycles} & \textbf{Fabricated} \\
\midrule
Within one era        & 617.2 d & 1.69 &  3.20\,\% \\
Across the join       & 988.9 d & 2.71 & 14.52\,\% \\
\midrule
\textbf{Worth of the join} & \textbf{+371.7 d} & \textbf{+1.02} & --- \\
\bottomrule
\end{tabular}
\end{center}

\panel{The join buys a year, and the gain is not free of doubt}{%
Accepting the estimated offset adds \textbf{371.7 days} and takes the record
past \textbf{2.7 annual cycles} --- enough, for the first time anywhere in this
project, to identify an annual component from a single continuous window.

\medskip
It also inherits the unresolved finding of Section~\ref{sec:changeover}: the
diurnal amplitude differs by 34\,\% across the changeover, so the two
instruments do not share a scale. A stretch spanning them is continuous in
level and inhomogeneous in gain.

\medskip
\textbf{The recommendation is therefore conditional.} For work in the difference
domain, or for anything where a change of scale partway through is tolerable,
take the 988.9-day stretch. For anything fitting one amplitude across the whole
window --- a harmonic decomposition, a seasonal model, an amplitude-based
diagnostic --- take the 617.2-day stretch inside the legacy era, which is
homogeneous and costs 3.2\,\% fabrication rather than 14.5\,\%.}

\begin{figure}[H]
\includegraphics{U_F09_recommended_stretch}
\caption{The recommended stretch with its provenance. Upper panel: the series,
with imputed values marked and their calibrated 90\,\% interval shaded. Lower
panel: where the imputed values are.}
\end{figure}

% ===========================================================================
\section{Verdicts}
\label{sec:verdicts}

\begin{center}
\small
\begin{tabularx}{\textwidth}{L{4.2cm}Y}
\toprule
\textbf{Question} & \textbf{Answer} \\
\midrule
how many stations, and for how long? &
\texttt{st01} 1119 days, last reading 2022-06-01; \texttt{st02} 2061 days, still
recording; \texttt{st03} 984 days, last reading 2023-06-24. The network failed
in stages, not at the changeover. \\
how complete is the target station? &
2061 of 2940 days, 70.1\,\%. 569 gaps holding 21\,098 hours. \\
where are the missing hours? &
\textbf{84.4\,\% in six gaps over 30 days.} The other 563 gaps together hold
15.6\,\%, and the 560 shorter than three days hold 8.9\,\%. \\
what is the missingness mechanism? &
\textbf{Outage-driven (MNAR).} Every covariate is missing on 100\,\% of the
hours the target is missing. No covariate-conditioned method can fill a real gap
here. \\
best filler, short gaps &
1 h: kalman, 1.04\,mdeg. 6 h: interpolation, 2.79. Below six hours the choice
barely matters. \\
best filler, long gaps &
24 h and beyond: seasonal-naive or the Kalman smoother, 4--14\,mdeg.
\textbf{Interpolation collapses past 12 h} and must not be used there. \\
the driver filler &
Lowest error at eight of eleven lengths, and \fail\ --- it cannot run on a real
gap. Measured, then refused. Refusing it costs 1.05\,mdeg on average. \\
uncertainty on a filled value &
Calibrated without distributional assumption; half-widths of 2--31\,mdeg
depending on method and length, consistently 2--3$\times$ the mean error. \\
longest stretch, nothing filled &
\textbf{108.2 days} (0.30 cycles), 2025-05-05 to 2025-08-21. \\
longest stretch, after imputation &
\textbf{988.9 days} (2.71 cycles), 2023-06-21 to 2026-03-06, 14.5\,\% invented
--- but it crosses the instrument changeover. \\
longest homogeneous stretch &
\textbf{617.2 days} (1.69 cycles), 2020-07-31 to 2022-04-09, 3.2\,\% invented,
one instrument throughout. \\
what still cannot be filled &
The 438-day, 104-day and 59-day outages. Raising the tolerance from 15 to
40\,mdeg buys \textbf{not one additional day}. \\
\bottomrule
\end{tabularx}
\end{center}

\panel{The recommendation}{%
\textbf{For work in the difference domain, or tolerant of a scale change
partway: the 988.9-day stretch}, June 2023 to March 2026, 14.5\,\% imputed, with
the per-value intervals attached. It is the only window in this project reaching
2.7 annual cycles.

\medskip
\textbf{For anything fitting one amplitude across the window: the 617.2-day
stretch}, July 2020 to April 2022, 3.2\,\% imputed, one instrument throughout.
It is shorter and it is homogeneous, and for a harmonic or seasonal model that
matters more than its length.

\medskip
In both cases the filled values are marked, and their uncertainty travels with
them.}

% ===========================================================================
\section{Caveats}

\textbf{The compensation coefficient is uncalibrated.} Every absolute level in
this dataset depends on a value never measured for these instruments. Relative
statements are far more robust than levels, and the first difference is
invariant to it entirely.

\textbf{The conformal intervals are calibrated on injected gaps.} Their
guarantee is marginal and holds under exchangeability between injected and real
gaps. Section~\ref{sec:missing} shows that assumption failing in one identified
respect --- real gaps have no drivers --- so the intervals are a floor on the
uncertainty, not a complete account of it.

\textbf{The joined series reconciles levels, not scales.} The 34\,\% difference
in diurnal amplitude across the changeover is unexplained and is not removed by
the offset.

\textbf{One station is analysed.} All three are built and classified, and only
\texttt{st02} is carried through the imputation work, by instruction and for
comparability with the two prediction studies.

% ===========================================================================
\section{Artefact inventory}

\begin{center}
\small
\begin{tabularx}{\textwidth}{L{3.4cm}Y}
\toprule
\textbf{Artefact} & \textbf{Content} \\
\midrule
\texttt{U\_01} -- \texttt{U\_03} & Station classification, era profiles, and the
changeover estimate. Written by the ingest \\
\texttt{U\_04} -- \texttt{U\_07} & Classification, era profile, monthly and
calendar completeness \\
\texttt{U\_08} -- \texttt{U\_10} & Every gap, the length bands, and the nine
major outages \\
\texttt{U\_11}, \texttt{U\_12} & Missingness mechanism, whole record and current
era \\
\texttt{U\_13} -- \texttt{U\_15} & Filler benchmark, conformal bands, and the
driver illusion \\
\texttt{U\_15b} & What excluding the unavailable filler costs, band by band \\
\texttt{U\_16} & The policies, one per tolerance \\
\texttt{U\_17}, \texttt{U\_18} & Stretch frontier: across the era join, and
confined within one era \\
\texttt{U\_19}, \texttt{U\_20} & The recommended stretch, and the verdicts \\
\addlinespace
\texttt{U\_F01}, \texttt{U\_F02} & Raw series per station; the era join \\
\texttt{U\_F03} -- \texttt{U\_F05} & Station completeness, channel availability,
calendar \\
\texttt{U\_F06} & Gap anatomy \\
\texttt{U\_F07} & Filler benchmark and calibrated intervals \\
\texttt{U\_F08}, \texttt{U\_F09} & Stretch frontier; the recommended stretch \\
\bottomrule
\end{tabularx}
\end{center}

Reproduction, from the study directory with \texttt{neuralprophet\_env} active:

\begin{quote}
\ttfamily\small
jupytext -{}-to ipynb build\_unified\_dataset.py \\
jupyter nbconvert -{}-to notebook -{}-execute -{}-inplace \\
\phantom{jupyter nbconvert }build\_unified\_dataset.ipynb \\[2mm]
jupytext -{}-to ipynb imputation\_study.py \\
jupyter nbconvert -{}-to notebook -{}-execute -{}-inplace \\
\phantom{jupyter nbconvert }imputation\_study.ipynb
\end{quote}

The ingest must run first: the study reads its output and nothing else.

\end{document}
